\subsection{DDT与LAT}
根据课程所学的差分分析和线性分析，这里简单复述一下S盒的差分分布表(DDT)和线性分布表(LAT)。

\subsubsection{DDT}
对于一个$m \times n$的S盒，即输入$m$比特、输出$n$比特的映射$S:\mathbb{F}_2^m \to \mathbb{F}_2^n$，记满足输入差分$\alpha$、输出差分$\beta$的输入集合为
$$
IN_S(\alpha ,\beta) = \{X \in \mathbb{F}_2^m :S(X) \oplus S(X \oplus \alpha) = \beta\}.
$$
DDT中的表项为该集合的元素个数，即
$$
\operatorname{DDT}[\alpha][\beta] = \# IN_S(\alpha,\beta).
$$
可见这是一张$2^m \times 2^n$的表，构建时可以先遍历输入差分$\alpha \in \mathbb{F}_2^m$共$2^m$种可能，内层循环遍历输入$X \in \mathbb{F}_2^m$共$2^m$种可能，根据计算出的输出差分更新表项。因此，构建DDT需要检查$2^m \times 2^m = 2^{2m}$组$(\alpha,X)$对应的差分传播。

\subsubsection{LAT}

对于同一个$m \times n$的S盒，LAT是一个$2^m \times 2^n$的表，其中第$\alpha$行第$\beta$列的表项为
$$
\operatorname{LAT}[\alpha][\beta]
=\#\{x\in\mathbb{F}_2^m\mid \alpha\cdot x\oplus\beta\cdot S(x)=0\}-2^{m-1}.
$$
这里$\cdot$表示二进制向量的模$2$内积。该表项是输入、输出掩码为$\alpha,\beta$时，使线性表达式$\alpha\cdot x\oplus\beta\cdot S(x)$等于$0$的输入个数与$2^{m-1}$的差，其值与输入总数$2^m$的比值恰好为线性近似的偏差：
$$
\epsilon(\alpha,\beta)
=\Pr[\alpha\cdot x=\beta\cdot S(x)]-\frac{1}{2}
=\frac{\operatorname{LAT}[\alpha][\beta]}{2^m}.
$$
其中概率是在输入$x$均匀分布于$\mathbb{F}_2^m$时计算的。构建该表时，可以先遍历输入掩码$\alpha\in\mathbb{F}_2^m$，内层遍历输出掩码$\beta\in\mathbb{F}_2^n$，最内层遍历输入$x\in\mathbb{F}_2^m$，统计线性表达式为$0$的输入个数。因此，共需要检查$2^m\times2^n\times2^m=2^{2m+n}$组$(\alpha,\beta,x)$对应的线性表达式。
